% frontmatter/czytam/mapa-del-curso.tex -- el año entero en una tabla,
% como frontmatter/lupa/mapa-del-curso.tex, en polaco: las palabras del
% texto de cada parte, medidas en polaco (content/czytam/progresion.json,
% redondeadas), y la última columna sin justificar, como en «Zdania».
\section*{Mapa roku}

\begin{center}
\renewcommand{\arraystretch}{1.4}
\begin{tabularx}{\linewidth}{@{}c c c c >{\raggedright\arraybackslash}X@{}}
\toprule
\textbf{Część} & \textbf{Dni} & \textbf{Pora roku} & \textbf{Słowa} & \textbf{Wielka tajemnica} \\
\midrule
1 & 1--65   & jesień & 45--75 & Przychodzi Hugo, powstaje Klub Lupy, a
w domu pojawiają się tajemnicze liściki podpisane \textit{Czytelnik X}.
Kto je pisze? \\
2 & 66--130 & zima & 75--90 & W Wigilię Lucía znajduje mapę skarbu,
którą babcia Rosa narysowała, kiedy miała osiem lat. \\
3 & 131--195 & wiosna & 90--105 & Dwadzieścia pięć lat temu jedna klasa
zakopała na szkolnym podwórku kapsułę czasu. Nikt nie pamięta gdzie. \\
4 & 196--260 & lato & 100--115 & Na wsi, z mapą babci: gdzie jest skarb
Rosy? \\
\bottomrule
\end{tabularx}
\end{center}

\vspace{6mm}
Każda część ma 13 tygodni po 5 dni, także w czasie ferii i wakacji.
Przez cały rok rośnie nie tylko długość tekstu (kolumna
\textit{Słowa}), ale też długość zdań i liczba ich części, czas, w
którym opowiadana jest historia -- teraźniejszy jesienią, przeszły od
zimy -- i rodzaj tekstu: liściki, spisy, listy, przepisy,
instrukcje, wywiady, teksty popularnonaukowe, pamiętnik, wiersze i
artykuły z gazetki.
\newpage
